\chapter{GENEL BİLGİLER}
\label{ch:genelbilgiler}

Bu bölümde tez çalışmasının kavramsal altyapısını oluşturan konular ele alınmaktadır.
Önce yapay zekâ destekli BT operasyonları (AIOps) kavramı ve bulut platformlarındaki kapasite yönetimi sorunları sunulmaktadır.
Ardından log analizi ve anomali tespiti yöntemleri incelenmekte, makine öğrenmesi modelleri ile açıklanabilir yapay zekâ yaklaşımları tartışılmaktadır.
Son olarak büyük dil modellerinin operasyonel zekayla entegrasyonu ve literatürdeki açık noktalar değerlendirilerek tezin konumlandığı alan belirlenmektedir.


\section{Yapay Zekâ Destekli BT Operasyonları (AIOps)}

BT sistemlerinin ölçeği büyüdükçe operasyonel verilerin hacmi (log, metrik, olay akışları) insan kapasitesinin çok üzerine çıkmaktadır.
Bu gerçeklik, makine öğrenmesi ve büyük veri tekniklerini BT operasyonları süreçlerine entegre eden AIOps (Artificial Intelligence for IT Operations) yaklaşımını doğurmuştur \autocite{dang2019}.

Notaro ve arkadaşları, AIOps literatürünü kapsamlı biçimde inceleyerek başarısızlık yönetimi için kullanılan yöntemleri anomali tespiti, kök neden analizi ve olay korelasyonu başlıkları altında sınıflandırmıştır \autocite{notaro2021}.
Bu sınıflandırmaya göre proaktif kapasite yönetimi, anomali tespitinin özel bir formu olarak değerlendirilebilir: kaynak kullanımının belirli bir eşiği aşmadan önce örüntüden sapması tespit edilerek önleyici aksiyon alınır.

AIOps'un başarısı büyük ölçüde veri kalitesine ve farklı kaynaklardan gelen sinyallerin etkin biçimde birleştirilmesine bağlıdır.
Metrikler, kesintisiz sayısal sinyal sağlarken; loglar, hata türleri, uyarı desenleri ve servis sağlık durumu hakkında metriklerin içermediği bağlamsal bilgi sunar.
Bu iki kaynağın korelasyonu, tek kaynaklı yaklaşımlara kıyasla daha erken ve daha kesin uyarı üretilmesine olanak tanımaktadır.


\section{Bulut Altyapılarında Kapasite Yönetimi}

Kapasite yönetimi, mevcut altyapının talepleri karşılamaya devam edip etmeyeceğini öngörme ve buna göre planlama yapma sürecidir.
Bulut ortamlarında bu süreç; vCPU, bellek ve disk gibi temel kaynakların kullanım trendlerinin izlenmesini ve gelecekteki tüketimin tahmin edilmesini kapsamaktadır.

Geleneksel yaklaşımlar iş yükü tahmininde ARIMA ve üstel düzeltme gibi istatistiksel zaman serisi yöntemlerine dayanmaktadır \autocite{calheiros2015}.
Bu yöntemler durağan ve periyodik örüntülü verilerde başarılı olurken, staj dönemleri ve anlık yük artışları gibi düzensiz olayları içeren bulut iş yüklerinde yetersiz kalabilmektedir.

Makine öğrenmesi yöntemleri bu sınırlılığı aşmak için önerilmiştir.
Özellik mühendisliği yoluyla geçmiş pencere istatistikleri, gecikme değerleri ve takvim bilgileri (iş günü, saat) modele dahil edilerek zaman serisi tahmini bir regresyon problemine dönüştürülmektedir.
Bu yaklaşımda her gözlem noktası, geçmiş zaman penceresinden türetilen özellik vektörüyle temsil edilmekte ve model bu vektörden gelecekteki kaynak kullanımını tahmin etmektedir.

OpenStack tabanlı özel bulut altyapılarında Nova (işlem), Cinder (depolama), Neutron (ağ) ve Keystone (kimlik) bileşenleri, Prometheus/Thanos ile saatlik metrik akışı sağlayarak makine öğrenmesi destekli kapasite yönetiminin veri altyapısını oluşturmaktadır \autocite{openstack2024}.


\section{Log Analizi ve Anomali Tespiti}

Log verileri, sistem olaylarının zaman damgalı metin kaydıdır.
Ham loglardan kapasite yönetimine anlamlı girdi oluşturmak birkaç aşamalı bir işlemi gerektirir: ayrıştırma ve sınıflandırma, zamansal agregasyon, anomali tespiti ve ön-arıza örüntü çıkarımı.

He ve arkadaşları, log ayrıştırma algoritmalarını kapsamlı biçimde değerlendirerek Drain, Spell ve AEL gibi yöntemlerin gerçek dünya log verisi üzerindeki etkinliğini karşılaştırmıştır \autocite{he2016}.
Log analizi için otomasyon yaklaşımlarını inceleyen daha kapsamlı bir tarama çalışması, ayrıştırma, kümeleme, sınıflandırma ve anomali tespiti adımlarını birleştiren bütünleşik pipeline yaklaşımlarının etkinliğini ortaya koymaktadır \autocite{he2021}.

Anomali tespiti açısından Chandola ve arkadaşları, noktasal, bağlamsal ve kolektif anomaliler arasındaki ayrımı tanımlayarak istatistiksel, kümeleme tabanlı ve sınıflandırma tabanlı yöntemleri sistematik biçimde sınıflandırmıştır \autocite{chandola2009}.
Derin öğrenme tabanlı yaklaşımlar ise özellikle yüksek boyutlu ve yapılandırılmamış veriler için güçlü alternatifler sunmaktadır \autocite{pang2021}.

Log dizisi anomali tespitinde önemli bir dönüm noktası olan DeepLog, sistem loglarını LSTM modeli ile modelleyerek normal log dizilerinden sapmaları tespit etmekte ve tanılamaktadır \autocite{du2017}.
Log verisi için derin öğrenme yaklaşımlarını derleyen güncel bir tarama çalışması, denetimsiz ve yarı denetimli yöntemlerin üretim ortamlarındaki kararsız log formatlarına karşı daha dayanıklı olduğunu göstermektedir \autocite{landauer2023}.

Bulut platformlarında loglar yalnızca hata bildirmez; arızadan saatler önce ortaya çıkan \textit{yavaş yanıt}, \textit{bellek baskısı} ve \textit{bağlantı havuzu uyarısı} gibi erken uyarı sinyalleri de içerir.
Bu sinyallerin zamansal örüntülerinin çıkarılması ve metrik verisiyle korelasyonu, proaktif izlemenin temel katkısını oluşturmaktadır.


\section{Makine Öğrenmesi ile Zaman Serisi Regresyonu}

Kapasite tahmini bu çalışmada bir regresyon problemi olarak ele alınmaktadır: geçmiş özellik vektörü verildiğinde $h$ saat sonraki kaynak kullanım oranı $\hat{y}$ tahmin edilmektedir.

\textbf{Doğrusal Regresyon (LR):}
Özellik vektörünün doğrusal kombinasyonuyla hedef değişkeni modelleyen temel yöntemdir.
Yorumlanabilirliği yüksektir; ancak doğrusal olmayan ilişkileri yakalayamaz.
Bu çalışmada StandardScaler ile normalizasyon uygulanmaktadır.

\textbf{Rastgele Orman (RF):}
Breiman tarafından önerilen topluluk öğrenmesi yöntemidir \autocite{breiman2001}.
Birbirinden bağımsız karar ağaçlarının tahminlerini ortalaması alarak oluşturur; aşırı öğrenmeye karşı dayanıklıdır ve özellik önem skoru hesaplamaya doğrudan olanak tanır.

\textbf{Gradyan Artırmalı Ağaçlar (GBT):}
Friedman tarafından geliştirilen ve artık (residual) tabanlı sıralı ağaç inşa eden yöntemdir \autocite{friedman2001}.
XGBoost \autocite{chen2016}, LightGBM \autocite{ke2017} ve CatBoost \autocite{prokhorenkova2018} bu temel yöntemin hesaplama verimliliğini ve doğruluğunu artıran modern uygulamalarıdır.

\textbf{Çok Katmanlı Algılayıcı (MLP):}
İleri beslemeli yapay sinir ağıdır.
Doğrusal olmayan ilişkileri yakalama kapasitesi yüksektir; ancak hiperparametre ayarı ve eğitim süresi diğer yöntemlere kıyasla daha fazla dikkat gerektirir.
Bu çalışmada StandardScaler ile normalizasyon uygulanmaktadır.

Model performansı dört metrikle değerlendirilmektedir: MAE (Ortalama Mutlak Hata), RMSE (Kök Ortalama Kare Hata), R\textsuperscript{2} (Belirleme Katsayısı) ve MAPE (Ortalama Mutlak Yüzdelik Hata).
Bu metrikler hem mutlak hata büyüklüğünü (MAE, RMSE) hem de varyansın ne kadarının model tarafından açıklandığını (R\textsuperscript{2}) ve yüzdesel sapmayı (MAPE) birlikte değerlendirmeye olanak tanımaktadır.


\section{Açıklanabilir Yapay Zekâ ve SHAP}

Makine öğrenmesi modellerinin karar destek süreçlerinde kullanılabilmesi için operatörlerin model çıktılarını anlayabilmesi kritik önem taşımaktadır.
Açıklanabilir yapay zekâ (XAI), kara kutu modellerin tahminlerini yorumlanabilir kılmayı hedeflemektedir.

SHAP (SHapley Additive exPlanations), oyun teorisindeki Shapley değerlerinden türetilen tutarlı bir özellik atıf yöntemidir \autocite{lundberg2017}.
Her özelliğin modelin tahmine katkısı, o özelliğin dahil edilip edilmediği tüm alt kümelerin kombinasyonları üzerinden hesaplanmakta; bu sayede hem yerel (tek tahmin) hem de global (tüm veri seti) açıklamalar üretilebilmektedir.

Ağaç tabanlı modeller için önerilen TreeSHAP algoritması, polinom zamanlı naif SHAP hesabını üstel karmaşıklıktan polinom karmaşıklığa indirerek büyük ölçekli modellerde uygulanabilir kılmaktadır \autocite{lundberg2020}.
Doğrusal modeller için ise LinearExplainer, katsayıların beklenti değerinden sapmasını doğrudan hesaplar; model-agnostik Kernel SHAP ise MLP gibi kara kutu modellerle kullanılır.

Bu çalışmada SHAP analizi üç amaçla kullanılmaktadır:
(i) hangi metrik ve log özelliklerinin kapasite tahminini en çok etkilediğinin belirlenmesi,
(ii) log sayımlarının metrik özelliklerle rekabet eden tahmin katkısının ölçülmesi ve
(iii) operatörlere üretilen önerilerin dayandığı kanıtların aktarılması.


\section{Büyük Dil Modelleri ve Operasyonel Zekâ}

Büyük dil modelleri (LLM), milyarlarca parametreli dönüştürücü (transformer) mimarisiyle eğitilen ve doğal dil anlama ile üretme kapasitesine sahip modellerdir.
GPT-3 \autocite{brown2020} ile birlikte bağlam içi öğrenme (in-context learning) kapasitesi gösterilmiş; GPT-4 \autocite{openai2023} ve Claude \autocite{anthropic2024} ile bu kapasite üretim ortamlarında kullanılabilir düzeye ulaşmıştır.

IT operasyonları alanında LLM'ler, yapılandırılmış sistem verilerini (metrik, log, uyarı) doğal dil önerilerine dönüştürme rolüyle kullanılmaktadır.
Bu yaklaşımın başarısı büyük ölçüde modele sağlanan bağlamın kalitesine bağlıdır: yalnızca ham metrik değerlerinin iletilmesi yerine, SHAP özellik önem sıralaması, log anomali sinyalleri, tahmin ufku ve servis sağlık durumu gibi bağlamsal bilgilerin eklenmesi daha spesifik ve uygulanabilir öneriler üretilmesini sağlamaktadır.

Bu çalışmada Claude Sonnet ve GPT-4o modelleri paralel olarak kullanılmakta; her iki modelin yanıtı tek bir operasyonel raporda birleştirilmektedir.
Modellere iletilen istem (prompt), kapasite durumu, log özeti, SHAP sıralaması, ön-arıza örüntüleri ve cascade risk skoru içeren yapılandırılmış bir bağlam bloğunu kapsamaktadır.


\section{Literatürdeki Boşluk ve Tezin Konumlandığı Alan}

İncelenen çalışmalar değerlendirildiğinde şu eksiklikler göze çarpmaktadır:

\begin{maddelendir}[label=\arabic*.]
    \item Mevcut AIOps çalışmalarının büyük bölümü log analizi ya da metrik bazlı kapasite tahminini ayrı ayrı ele almaktadır.
    Her iki veri kaynağını tek bir özellik uzayında birleştiren ve korelasyonun tahmini nasıl iyileştirdiğini nicel olarak ölçen uygulamalı çalışmalar sınırlı kalmaktadır.

    \item Ön-arıza örüntü çıkarımı (pre-failure pattern mining) kapasite tahminiyle birlikte ele alınmamaktadır.
    Geçmiş incident'lerin öncesindeki log imzalarının sistematik biçimde çıkarılıp gelecekteki uyarı üretimine katkı sağlaması literatürde yeterince araştırılmamıştır.

    \item SHAP analizinin kapasite yönetimi bağlamında uygulanması ve log özelliklerinin tahmin katkısının metrik özellikleriyle karşılaştırılması spesifik olarak ele alınmamıştır.

    \item LLM entegrasyonu genellikle log-to-text veya alert-to-text dönüşümlerine odaklanmaktadır.
    Metrik, log, SHAP ve cascade bilgisini birleştiren zenginleştirilmiş bağlam ile proaktif öneri üretimi yeterince araştırılmamıştır.
\end{maddelendir}

Bu tez, söz konusu boşlukları hedef alarak OpenStack tabanlı gerçek bir bulut altyapısı üzerinde log--metrik korelasyonuna dayalı, açıklanabilir ve LLM destekli proaktif bir kapasite yönetim sistemi geliştirmektedir.
